% BHU PYQ -- Geography: Physical Geography (2024-25 NEP)
\documentclass[11pt,a4paper]{article}
\usepackage[a4paper,top=2.5cm,bottom=2.5cm,left=2.8cm,right=2.8cm,headheight=14pt]{geometry}
\usepackage{amsmath,amssymb}
\usepackage{fontspec}
\setmainfont{Kohinoor Devanagari}
\usepackage{microtype,fancyhdr,enumitem,hyperref,lastpage,parskip}

\hypersetup{colorlinks=true,urlcolor=black,linkcolor=black,pdftitle={GGRMJ11: Physical Geography -- BHU NEP 2024-25}}

\pagestyle{fancy}\fancyhf{}
\renewcommand{\headrulewidth}{0.4pt}\renewcommand{\footrulewidth}{0.4pt}
\fancyhead[L]{\small \textit{Banaras Hindu University}}
\fancyhead[R]{\small \textit{GGRMJ11: Physical Geography (2024-25)}}
\fancyfoot[C]{\small Page \thepage\ of \pageref{LastPage}}

\newlist{parts}{enumerate}{2}
\setlist[parts,1]{label=(\alph*),leftmargin=*,itemsep=4pt,topsep=3pt}
\setlist[parts,2]{label=(\roman*),leftmargin=*,itemsep=2pt,topsep=2pt}

\newcommand{\pts}[1]{\hfill \textit{[#1]}}

\begin{document}

\begin{center}
  {\large\bfseries BANARAS HINDU UNIVERSITY}\\[2pt]
  {\normalsize B. A. / B. Sc. (Hons.) Semester I Examination, 2024-25 (NEP)}\\[6pt]
  \rule{0.8\linewidth}{0.5pt}\\[6pt]
  {\large\bfseries GEOGRAPHY}\\[4pt]
  {\large\bfseries Paper Code: GGRMJ-11 / GGRMN-11 : Physical Geography}\\[4pt]
  {\normalsize Time: 2:30 Hours \hfill Full Marks: 70}\\[2pt]
  {\small REF NO. 174}
\end{center}
\rule{\linewidth}{0.4pt}
\smallskip

\noindent \textit{Instructions: Answer five questions in all. All questions carry equal marks (14 marks each). Question No.\ 1 is compulsory. Answer each part of Question No.\ 1 in approximately 200 words and remaining questions in approximately 1000 words. Illustrate your answer with suitable maps and diagrams. Write your Roll No.\ at the top immediately on receipt of this question paper.}

\smallskip
\noindent \textit{कुल पाँच प्रश्नों के उत्तर दीजिए। सभी प्रश्नों के अंक समान हैं। प्रश्न सं० 1 अनिवार्य है। प्रश्न संख्या 1 के प्रत्येक भाग का उत्तर लगभग 200 शब्दों में एवं शेष प्रश्नों का उत्तर लगभग 1000 शब्दों में दीजिए। अपने उत्तर की पुष्टि उपयुक्त मानचित्रों तथा आरेखों द्वारा कीजिए।}

\bigskip

\begin{enumerate}[label=\textbf{\arabic*.},leftmargin=*,itemsep=14pt]

  \item Write short notes on the following:\pts{4 $\times$ 3.5 = 14}
  
  निम्नलिखित पर संक्षिप्त टिप्पणियाँ लिखिए:
  \begin{parts}
    \item Crust / भू-पर्पटी
    \item Local winds / स्थानीय पवनें
    \item Coral reefs / प्रवाल भित्तियाँ
    \item Metamorphic rock / कायान्तरित शैल
  \end{parts}

  \item Describe the works of running water with suitable examples and diagrams.\pts{14}
  
  प्रवाहित जल (सरिता) के कार्यों का उपयुक्त उदाहरणों एवं आरेखों की सहायता से वर्णन कीजिए।

  \item Explain any one hypothesis of the origin of the Earth.\pts{14}
  
  पृथ्वी की उत्पत्ति की किसी एक परिकल्पना की व्याख्या कीजिए।

  \item Describe the composition and structure of the atmosphere.\pts{14}
  
  वायुमण्डल के संगठन एवं संरचना का वर्णन कीजिए।

  \item Describe the air pressure belts. How do they affect the wind circulation?\pts{14}
  
  वायुदाब पेटियों का वर्णन कीजिए। ये वायु संचार को किस प्रकार प्रभावित करती हैं?

  \item Describe the ocean currents of the Atlantic Ocean.\pts{14}
  
  अन्धा महासागर (अटलांटिक महासागर) की महासागरीय धाराओं का वर्णन कीजिए।

  \item Explain the temperature distribution of the oceans.\pts{14}
  
  महासागरों के तापमान वितरण की व्याख्या कीजिए।

\end{enumerate}

\end{document}
